\clearpage\thispagestyle{otFront}\phantomsection\label{mode-emploi}\pdfbookmark[0]{Mode d’emploi}{bm-mode-emploi}\otPageTitre{Mode d’emploi}{Des repères stables pour travailler seul}\fontsize{11}{15}\selectfont

\begin{tabularx}{\linewidth}{@{}lY@{}}\toprule
Une pastille & Découvrir : calcul direct, vocabulaire, lecture d’une figure ou première application.\\[3mm]
Deux pastilles & S’entraîner : choisir un outil, combiner des étapes et vérifier.\\[3mm]
Trois pastilles & Résoudre : organiser une recherche, justifier, interpréter et contrôler.\\\bottomrule
\end{tabularx}\par\medskip
Ces niveaux sont des indications sur les tâches, pas des étiquettes sur les
élèves. Revenir à une question simple peut aider à débloquer un problème.\par\medskip
\begin{exemple}[Lire une correction]Pour $3(x-2)=x+8$, on développe : $3x-6=x+8$. On obtient $2x=14$, donc $x=7$. Le contrôle donne $3(7-2)=15=7+8$. Le résultat est utile parce que chaque étape peut être expliquée et refaite.\end{exemple}
\begin{attention}Lire les unités, les hypothèses, les mots « strictement », « au moins » et l’arrondi demandé. Une figure ressemblant à un triangle rectangle ne prouve pas qu’il l’est. Un modèle de hasard doit préciser les chances des issues.\end{attention}
Les liens du sommaire et des corrections permettent de naviguer. Pour les
scripts, les indentations et les lignes « fin » délimitent les blocs.
Après une correction, fermer le livre et refaire le passage décisif.

\clearpage
